\documentclass[11pt]{amsart}
\usepackage{amsmath,amssymb,amsthm}
\usepackage[margin=1.1in]{geometry}
\usepackage{newtxtext,newtxmath}
\usepackage{graphicx}
\usepackage{microtype}
\usepackage[colorlinks=true, linkcolor=blue, citecolor=blue, urlcolor=blue]{hyperref}
\usepackage{listings}
\usepackage{upquote}
\lstdefinelanguage{Lean}{
  morekeywords={def,theorem,open,structure,where,fun},
  sensitive=true,
  morecomment=[l]{--},
}
\lstset{language=Lean, basicstyle=\small\fontfamily{lmtt}\selectfont\upshape, keywordstyle=\bfseries,
  columns=fullflexible, keepspaces=true, breaklines=false,
  xleftmargin=1em, aboveskip=6pt, belowskip=6pt,
  literate={ℕ}{{$\mathbb{N}$}}1 {ℝ}{{$\mathbb{R}$}}1 {ℤ}{{$\mathbb{Z}$}}1
    {→}{{$\rightarrow$}}1 {×}{{$\times$}}1 {∃}{{$\exists$}}1 {∀}{{$\forall$}}1
    {ε}{{$\varepsilon$}}1 {≤}{{$\le$}}1 {≥}{{$\ge$}}1 {∧}{{$\wedge$}}1
    {¬}{{$\neg$}}1 {∣}{{$\mid$}}1 {∈}{{$\in$}}1 {₀}{{$_0$}}1 {Γ}{{$\Gamma$}}1
    {ω}{{$\omega$}}1 {κ}{{$\kappa$}}1 {∑}{{$\sum$}}1
}

\newtheorem{lemma}{Lemma}[section]
\newtheorem{remark}[lemma]{Remark}
\newtheorem*{theoremA}{Theorem}

\setlength{\emergencystretch}{3em}

\newcommand{\F}{\mathbb{F}}
\newcommand{\N}{\mathbb{N}}
\newcommand{\size}{\operatorname{size}}
\newcommand{\DFT}{\operatorname{DFT}}

\begin{document}

\title[The prime-field DFT below \texorpdfstring{$N\log N$}{N log N}]{The discrete Fourier transform over prime fields\\ below \texorpdfstring{$N \log N$}{N log N} on multitape machines}

\author{Dan Brown$^*$}
\email{jdb1729@gmail.com}
\thanks{$^*$Author of \emph{The Da Vinci Code}.}

\date{October 17, 2026}

\begin{abstract}
A fixed finite multi-stack Turing machine, in the model \texttt{Turing.FinTM2} of Lean's
Mathlib, computes the discrete Fourier transform of length $n$ over a prime field
$\F_p$, given $p$, a primitive $n$-th root of unity and the $n$ entries in binary of fixed
width, in worst-case time $O\bigl(N (\log N)^{1-\kappa}\bigr)$, where $N$ is the input length
and $\kappa > 0$ is the saving of the formally verified integer multiplier below $n \log n$.
The reduction is the standard reduction of the transform to one integer product: Bluestein's
chirp, written with binomial coefficients so that no square root of the root of unity is
needed, Kronecker packing, and Barrett reduction with a reciprocal computed by Newton's
iteration. The contribution is the explicit sub-$N\log N$ bound on multitape machines,
inherited from the multiplier, and a statement of it in Lean that uses Mathlib definitions
only. The Lean proof is in progress; we list what it proves.
\end{abstract}

\maketitle

\section{Introduction}

For a prime $p$, an integer $n \ge 1$, a primitive $n$-th root of unity $\omega \in \F_p$ and a
vector $a \in \F_p^n$, the discrete Fourier transform is
\[
  \DFT_\omega(a)_k = \sum_{j=0}^{n-1} a_j\, \omega^{jk} \qquad (0 \le k < n).
\]
Over a field whose transform lengths are smooth, the fast Fourier transform computes it with
$O(n \log n)$ field operations. The bit complexity on a multitape Turing machine is a different
question: the field elements have $w = \size(p)$ bits, $n$ need not be smooth, and the machine
must move its data along tapes. The classical route reduces the transform to a single
integer multiplication of $O(N)$ bits, where $N \asymp n w$ is the input length, so the
transform costs the time of that product plus $O(N)$ for the surrounding passes and
$O(n)$ multiplications of $w$-bit numbers.

The time of integer multiplication on multitape machines has been bounded by
$O(n \log n)$ by Harvey and van der Hoeven \cite{HvdH}, and below $n \log n$ by the
OpenAI manuscript \cite{OAI}: time $O\bigl(n (\log n)^{1-\kappa}\bigr)$ for a positive
$\kappa$. A complete formal proof of a bound of this form, for one fixed machine in Mathlib's
multi-stack model, is the Lean development \texttt{IntegerMultBounds} \cite{IMB}; its theorem
\texttt{MathlibTM2.mulInTime} gives the saving $\kappa = 88279\cdot 2^{-46} \approx 2^{-29.57}$.
The reduction then gives the same saving for the transform.

\begin{theoremA}
Let $\kappa > 0$ be such that some fixed finite multi-stack machine multiplies natural numbers
in time $O\bigl(n (\log n)^{1-\kappa}\bigr)$, with binary inputs and outputs. Then a fixed
finite multi-stack machine computes $\DFT_\omega(a)$ for every prime $p$, every primitive
$n$-th root of unity $\omega \in \F_p$ and every $a \in \F_p^n$, with the input written as $p$,
$\omega$, $a_0, \dots, a_{n-1}$ in exactly $\size(p)$ binary digits each, in time
\[
  O\bigl(N (\log N)^{1-\kappa}\bigr),
\]
where $N = (n+2)(\size(p)+1)$ is the input length. The implicit constant depends neither on $p$
nor on $n$. In particular this holds with $\kappa = 88279\cdot 2^{-46}$.
\end{theoremA}

The bound is uniform in the relation between $n$ and $p$: it covers $n = 2$ with $p$ of
millions of bits as well as $n = p - 1$. The first regime requires that every one-time setup
on $p$ costs $O(M(w))$, the time of one product of $w$-bit numbers, without a logarithmic
number of such products; Section~\ref{sec:cost} explains where this forces the choice of the
modular reduction.

Section~\ref{sec:statement} gives the statement in Lean and explains why the inputs have
fixed width. Section~\ref{sec:reduction} gives the reduction, Section~\ref{sec:cost} the
cost accounting, and Section~\ref{sec:formal} the state of the formalization.

\section{The statement}
\label{sec:statement}

The Lean statement is \texttt{DFTStatement} in \texttt{MathlibTM2/DFT/Statement.lean}. Apart from
the few definitions shown, everything in it is Mathlib vocabulary: the alphabet $\Gamma'$, the
encoding \texttt{encodeNat} of $\N$ in binary, \texttt{ZMod}, \texttt{IsPrimitiveRoot},
\texttt{TM2ComputableInTime} and the asymptotic relation \texttt{=O[atTop]}. Numbers are
written least significant bit first, each followed by a comma.

\noindent\begin{minipage}{\linewidth}
\begin{lstlisting}
def encodeFixed (w x : ℕ) : List Γ' :=
  List.ofFn fun i : Fin w => Γ'.bit (x.testBit i)

def encodeFixedList (w : ℕ) (l : List ℕ) : List Γ' :=
  l.flatMap fun x => encodeFixed w x ++ [Γ'.comma]

def encodeNatList (l : List ℕ) : List Γ' :=
  l.flatMap fun x => (encodeNat x).map inclusionBoolΓ' ++ [Γ'.comma]

def dft {p n : ℕ} (ω : ZMod p) (a : Fin n → ZMod p) (k : Fin n) : ZMod p :=
  ∑ j, a j * ω ^ ((j : ℕ) * k)

structure DFTInstance where
  p : ℕ
  [prime : Fact p.Prime]
  n : ℕ
  ω : ZMod p
  primitive : IsPrimitiveRoot ω n
  a : Fin n → ZMod p

def DFTInstance.encode (x : DFTInstance) : List Γ' :=
  encodeFixedList (Nat.size x.p) (x.p :: x.ω.val :: List.ofFn fun j => (x.a j).val)

def DFTInstance.transform (x : DFTInstance) : List ℕ :=
  List.ofFn fun k => (dft x.ω x.a k).val

def DFTInTime (κ : ℝ) : Prop :=
  ∃ M : TM2ComputableInTime DFTInstance.encode encodeNatList DFTInstance.transform,
    (fun N : ℕ => (M.time N : ℝ)) =O[atTop]
      fun N : ℕ => (N : ℝ) * Real.log N ^ (1 - κ)

def DFTStatement : Prop := ∃ κ : ℝ, 0 < κ ∧ DFTInTime κ
\end{lstlisting}
\end{minipage}

\texttt{TM2ComputableInTime} is Mathlib's notion: a machine of type \texttt{Turing.FinTM2}, a
time function of the input length, and a proof that on every input the machine halts with the
encoded output within that time. The time function therefore bounds the worst case over all
instances of a given length. The output is the list of canonical representatives
$\DFT_\omega(a)_k \in \{0, \dots, p-1\}$ in shortest binary form; its length is at most
$n (\size(p) + 1) \le N$.

\begin{remark}[Fixed-width inputs]
\label{rem:width}
With shortest binary forms on input as well, no bound $O\bigl(N (\log N)^{1-\kappa}\bigr)$ is
possible. Take $n = p - 1$, a generator $\omega$ of $\F_p^\times$, and
$a = (0, 1, 0, \dots, 0)$. In shortest form each zero entry is a single comma, so the input has about
$p + 2\log_2 p$ symbols. The transform is $\DFT_\omega(a)_k = \omega^k$, which runs over all of $\F_p^\times$, so
the output has $\sum_{x=1}^{p-1} (\size(x) + 1) \asymp p \log p$ symbols. A machine writes at
most one output symbol per step, so its time on these inputs is at least of order $N \log N$.
With every entry written in exactly $\size(p)$ digits, the output is never longer than the
input, and this lower bound disappears.
\end{remark}

\section{The reduction}
\label{sec:reduction}

Fix an instance with $n \ge 2$ and write $w = \size(p)$, so $2^{w-1} \le p < 2^w$. Since
$\omega$ has order $n$ in $\F_p^\times$, $n \mid p - 1$; hence $n < p$, $\size(n) \le w$, and
$p$ is odd. The cases $n \le 1$ are immediate: for $n = 0$ the output is empty, and for $n = 1$
we have $\omega = 1$ and $\DFT_\omega(a) = a$.

\subsection{The chirp without a square root}
Bluestein \cite{Blu} turns the transform into a correlation through the identity
$jk = \tfrac{1}{2}\bigl((j+k)^2 - j^2 - k^2\bigr)$, which uses a square root of $\omega$. The
binomial form avoids it:
\[
  jk = \binom{j+k}{2} - \binom{j}{2} - \binom{k}{2},
\]
an identity of integers. Put $\omega' = \omega^{n-1}$, so $\omega'\omega = 1$, and
$c_m = \omega^{\binom{m}{2}}$, $c'_m = \omega'^{\binom{m}{2}}$. Then
$\omega^{jk} = c'_j\, c'_k\, c_{j+k}$, and
\[
  \DFT_\omega(a)_k = c'_k \sum_{j=0}^{n-1} u_j\, c_{j+k}, \qquad u_j = a_j\, c'_j .
\]
The sum is a correlation of the vector $u$ of length $n$ with the chirp $c$ of length
$2n - 1$.

\subsection{Kronecker packing}
Let $\bar u_j, \bar c_m \in \{0, \dots, p-1\}$ be the representatives and $B = 2^s$ with
$s = 2w + \size(n)$, so that $n (p-1)^2 < B$. Pack the reversed vector and the chirp as integers,
\[
  U = \sum_{i=0}^{n-1} \bar u_{n-1-i}\, B^i, \qquad C = \sum_{m=0}^{2n-2} \bar c_m\, B^m .
\]
The coefficient of $B^{n-1+k}$ in the polynomial product is
$T_k = \sum_{j=0}^{n-1} \bar u_j\, \bar c_{j+k}$, and every coefficient of the product is a sum
of at most $n$ terms below $(p-1)^2$, hence below $B$. So no carries occur, and $T_k$ is the
digit of index $n - 1 + k$ of the integer $UC$ in base $B$:
\[
  T_k = \bigl\lfloor UC / B^{n-1+k} \bigr\rfloor \bmod B .
\]
The two factors have $ns$ and $(2n-1)s$ bits, so the single product has input length at most
$3ns + O(1) = O(N)$. Then $\DFT_\omega(a)_k = c'_k\, T_k$ in $\F_p$.

\subsection{Reduction modulo \texorpdfstring{$p$}{p}}
The remaining arithmetic is $O(n)$ products of residues and $n$ reductions of $T_k < n p^2$.
Both are Barrett reductions \cite{Bar}. With $\mu = \lfloor 2^{s'}/p \rfloor$ and $x < 2^{s'}$,
the quotient $q = \lfloor x\mu / 2^{s'} \rfloor$ satisfies
\[
  qp \le x < qp + 2p ,
\]
so $x \bmod p$ is $x - qp$ after at most two subtractions of $p$. One reduction costs two
products of $O(s')$-bit numbers. We use $s' = 2w$ for products of residues and $s' = s$ for the
sums $T_k$.

The constant $\mu$ is computed by Newton's iteration for the reciprocal with precision
doubling \cite[\S 4.3.3]{Knu}. For a divisor $d$ of $2k$ bits and $a$ the reciprocal of its top
$k$ bits at precision $k$, the step
\[
  y = 2^{k+1} a - \bigl(\lfloor d a^2 / 4^k \rfloor + 1\bigr)
\]
satisfies $y d < 16^k < (y + 9) d$, so the reciprocal $\lfloor 16^k / d \rfloor$ at precision
$2k$ is $y$ plus at most eight corrections. The levels have precisions $1, 2, 4, \dots$ up to
$O(w)$, each with three products, so the total is a geometric sum, $O(M(w))$.

\begin{remark}[Montgomery reduction]
\label{rem:montgomery}
Montgomery's reduction \cite{Mon} replaces division by $p$ with division by $R = 2^r$,
$R > np$, given the constant $-p^{-1} \bmod R$. That constant is computed by Hensel lifting
$y \mapsto y(2 + py)$ in $O(M(w))$ time, again a geometric sum. Entering Montgomery form,
however, needs $R^2 \bmod p$, a reduction of a $2r$-bit number modulo $p$. Computed by
$\log_2 r$ Montgomery squarings, it costs $O(M(w) \log w)$, which exceeds the budget of
Section~\ref{sec:cost} by a factor of order $\log w$ when $n$ is bounded. Computing it in $O(M(w))$ requires a division by $p$, and with the
reciprocal $\mu$ in hand Barrett reduction performs every reduction of the pipeline directly.
\end{remark}

\section{Cost}
\label{sec:cost}

Let $M(L)$ be the time of the multiplier on inputs of total length $L$, replaced by a monotone
majorant with $M(L) \le C_0\, L (\log(L+2))^{1-\kappa} + C_0$; then $M(L)/L$ is bounded by a
nondecreasing function, and $M(cL) = O(M(L))$ for fixed $c$. Every product of $O(w)$-bit numbers
costs $O(M(w))$ including the transfer of its operands and its result, and every pass over the
data costs $O(N)$. With $N = (n+2)(w+1)$ and $n \ge 2$:

\begin{enumerate}
\item reading the input and printing the output: $O(N)$;
\item the reciprocals $\lfloor 4^w / p \rfloor$ and $\lfloor 2^s / p \rfloor$ by Newton's
  iteration: $O(M(w))$;
\item $\omega' = \omega^{n-1}$ by repeated squaring: $O(\log n \cdot M(w))$;
\item the chirps $c_m$ $(m < 2n-1)$ and $c'_m$ $(m < n)$, from $c_{m+1} = c_m\, \omega^m$ and
  $\omega^{m+1} = \omega^m \omega$: $O(n\, M(w))$;
\item the chirped input $u_j = a_j c'_j$, packed in reverse: $O(n\, M(w) + N)$;
\item the product $UC$: $M(9N)$, since $3ns \le 9 n w \le 9N$;
\item the digits $T_k$, their reductions and the products $c'_k T_k$:
  $O\bigl(n\, M(w) + N\bigr)$, since $s \le 3w$.
\end{enumerate}

Since $nw \le N$ and $M(w)/w$ is bounded by $O\bigl((\log N)^{1-\kappa}\bigr)$, each term
$n\, M(w)$, $\log n \cdot M(w)$ and $M(w)$ is $O\bigl(N (\log N)^{1-\kappa}\bigr)$, and so is
$M(9N)$. The total is $O\bigl(N (\log N)^{1-\kappa}\bigr)$, uniformly in $n$ and $p$.

\begin{remark}[Tiny $n$, huge $p$]
For $n = 2$ the input length is $N = 4(w+1)$ and the budget is $O\bigl(w (\log w)^{1-\kappa}\bigr)$,
the time of a constant number of products of $w$-bit numbers. Every step of the list fits: the
Newton reciprocal is a geometric sum of products, and the other steps use $O(n)$ products.
A setup with $\Theta(\log w)$ products of $w$-bit numbers does not fit, which is the reason for
the choice in Remark~\ref{rem:montgomery}. For $n$ close to $p$ the product $UC$ dominates.
\end{remark}

\section{Formalization}
\label{sec:formal}

The statement of Section~\ref{sec:statement} is in \texttt{MathlibTM2/DFT/Statement.lean} of the
repository \cite{IMB}, which imports Mathlib only; its examples check \texttt{encodeFixedList},
\texttt{encodeNatList} and one transform over $\F_5$ by \texttt{decide}. The proof of
\texttt{DFTStatement} is in progress. The development takes the multiplier as a parameter,
a structure \texttt{HostMultiplier}~$\kappa$ holding a machine table, its time function and the
bound $O\bigl(n(\log n)^{1-\kappa}\bigr)$, so nothing in it depends on the value of $\kappa$.
The following parts are proved, with no \texttt{sorry}:

\begin{itemize}
\item \texttt{DFT/Algorithm.lean}: the arithmetic of Section~\ref{sec:reduction} on natural
  numbers. \texttt{chirp\_identity} is the binomial form of Bluestein's identity with
  $\omega'\omega = 1$; \texttt{pack\_mul}, \texttt{pack\_div\_pow\_mod}, \texttt{conv\_le} and
  \texttt{conv\_reflect} are Kronecker packing and digit read-back; \texttt{correlation\_cast}
  allows arbitrary representatives; \texttt{fastDFT\_transform} states that the whole
  computation, with reductions modulo $p$, equals \texttt{DFTInstance.transform} when
  $n(p-1)^2 < B$.
\item \texttt{DFT/Reduction.lean} and \texttt{DFT/Newton.lean}: Montgomery's REDC
  (\texttt{redc\_lt}, \texttt{redc\_cast}), the Hensel constant (\texttt{montInv\_dvd}), Barrett
  reduction (\texttt{barrett\_mod}, \texttt{barrett2\_le}, \texttt{barrett2\_lt}) and the Newton
  step (\texttt{newton\_step}).
\item \texttt{DFT/Keys.lean}: a machine with two $\Gamma'$ stacks around the host register file
  of \cite{IMB}, the transfer of host programs (\texttt{exec\_hostProg}), and the packaging into
  \texttt{TM2ComputableInTime} and \texttt{DFTInTime} (\texttt{computableInTime},
  \texttt{dftInTime\_of\_computes}).
\item \texttt{DFT/IO.lean}, \texttt{DFT/Fields.lean}, \texttt{DFT/Parse.lean}: reading the
  fixed-width input and printing the output (\texttt{parse\_exec}, \texttt{print\_exec}).
\item \texttt{DFT/MulIO.lean}, \texttt{DFT/MulCall.lean}: the multiplier as a subroutine on
  internal registers (\texttt{call\_exec}, \texttt{mulCall\_exec}).
\item \texttt{DFT/Trunc.lean}, \texttt{DFT/RedcTail.lean}, \texttt{DFT/Redc.lean},
  \texttt{DFT/Montgomery.lean}: REDC and Montgomery multiplication on the host
  (\texttt{redc\_exec}, \texttt{montMul\_exec}).
\item \texttt{DFT/Hensel.lean}, \texttt{DFT/HenselHost.lean}, \texttt{DFT/HenselLoop.lean}: the
  Hensel iteration with its step bound (\texttt{henselRound\_exec}, \texttt{hensel\_exec}).
\item \texttt{DFT/NewtonHost.lean}, \texttt{NewtonSetup.lean}, \texttt{NewtonBlocks.lean} and
  \texttt{NewtonGuess.lean}: the host blocks of one Newton level (\texttt{newtonPrep\_exec},
  \texttt{newtonLoad\_exec}, \texttt{newtonSq\_exec}, \texttt{newtonGuessProg\_exec},
  \texttt{newtonFinish\_exec}) and the correction loop (\texttt{corrLoop\_exec}).
\end{itemize}

Still to be written are the Newton iteration over all levels and Barrett reduction on the
host, the pipeline of Sections~\ref{sec:reduction} and~\ref{sec:cost} as one host procedure
with its step bound, the top level that turns it into \texttt{HostMultiplier}~$\kappa \to$
\texttt{DFTInTime}~$\kappa$, and the final theorem \texttt{dft\_in\_time : DFTStatement},
obtained from the multiplier of \texttt{MathlibTM2.mulInTime}, with its axiom audit.

\begin{thebibliography}{9}

\bibitem{Bar} P. Barrett,
\emph{Implementing the Rivest Shamir and Adleman public key encryption algorithm on a standard
digital signal processor}, Advances in Cryptology -- CRYPTO '86, Lecture Notes in Comput.
Sci. \textbf{263}, Springer, 1987, 311--323.

\bibitem{Blu} L. I. Bluestein,
\emph{A linear filtering approach to the computation of discrete Fourier transform},
IEEE Trans. Audio Electroacoust. \textbf{18} (1970), 451--455.

\bibitem{HvdH} D. Harvey, J. van der Hoeven,
\emph{Integer multiplication in time $O(n \log n)$}, Ann. of Math. \textbf{193} (2021),
563--617.

\bibitem{IMB} \texttt{IntegerMultBounds}: integer multiplication below $n \log n$ in Lean on
Mathlib's finite Turing machines,
\url{https://github.com/platypii/integer-mult-bounds-lean}.

\bibitem{Knu} D. E. Knuth,
\emph{The Art of Computer Programming}, Vol.~2: \emph{Seminumerical Algorithms}, 3rd ed.,
Addison-Wesley, 1997.

\bibitem{Mon} P. L. Montgomery,
\emph{Modular multiplication without trial division}, Math. Comp. \textbf{44} (1985),
519--521.

\bibitem{OAI} OpenAI,
\emph{Integer multiplication below $n \log n$}, manuscript, September 23, 2026,
\url{https://github.com/openai/math}.

\end{thebibliography}

\end{document}
